(1) 在子午光线近似中，纤芯光线与轴的夹角记为 β，碰到芯包侧壁时入射角相对于法线为 \(\pi/2-\beta\)。全内反射要求 \(\sin(\pi/2-\beta)=\cos\beta\ge n_2/n_1\)，因此 \(\sin\beta\le\sqrt{1-(n_2/n_1)^2}\)。在输入端 Snell 定律给 \(n_0\sin\theta_0=n_1\sin\beta\)，故
\[
n_0\sin\theta_{\max}=\sqrt{n_1^2-n_2^2}=:\mathrm{NA}.
\]
若 NA>n₀，物理接受角最多为 π/2；通常光纤满足 NA≤n₀。

(2) 令相对折射率差 \(\Delta=(n_1-n_2)/n_1\ll1\)。因 \(n_1^2-n_2^2=(n_1-n_2)(n_1+n_2)\simeq2n_1^2\Delta\)，得
\[
\mathrm{NA}\simeq n_1\sqrt{2\Delta}.
\]
因此即便芯包折射率差相对很小，也可以通过高数值孔径改变入射光的接受角，但该近似仅在弱导波条件成立。

(3) 波动光学定义归一化频率
\[
V=\frac{2\pi a}{\lambda}\sqrt{n_1^2-n_2^2}
=\frac{2\pi a\,\mathrm{NA}}{\lambda}.
\]
弱导波阶跃光纤中 LP11 模首次截止对应 \(V\simeq2.405\)（来自第一类贝塞尔函数的首个正零点），所以仅支持基模 LP01 的通常单模判据为 \(V<2.405\)，即
\[
\lambda>\frac{2\pi a\,\mathrm{NA}}{2.405}.
\]
基模本身无同样的正截止频率；实际双折射、弯曲与损耗会调整实用工作范围。

(4) 几何光学接受锥来自全内反射的射线条件，适于多模式、波长远小于纤芯直径的情况；单模截止由横向 Maxwell 本征问题及圆柱边界条件决定，并依赖无量纲参数 V。射线角限制不能替代模式数判据，因为导波模式是整个横截面上的相干场结构。
